% ============================================================
% Informe: Mandriladora portátil y sistema de recargue
% Cada sección vive en su carpeta (NN-nombre/nombre.txt) con su
% referencias.bib y sus figuras en recursos/. Compilar con:
%   pdflatex Mandriladora -> bibtex Mandriladora -> pdflatex x2
% ============================================================
\documentclass[12pt,a4paper]{article}

% ---------- Idioma y codificación ----------
\usepackage[utf8]{inputenc}
\usepackage[T1]{fontenc}
\usepackage[spanish]{babel}
\addto\captionsspanish{\renewcommand{\tablename}{Tabla}\renewcommand{\listtablename}{Índice de tablas}}

% ---------- Página ----------
\usepackage{geometry}
\geometry{a4paper, margin=2.5cm}

% ---------- Matemáticas, figuras y tablas ----------
\usepackage{amsmath,amssymb,amsfonts}
\usepackage{graphicx}
\usepackage{subcaption}
\usepackage{float}
\usepackage{booktabs}
\usepackage{tabularx}
\usepackage{array}
\usepackage{textcomp}
\usepackage{xcolor}
\usepackage{enumitem}
\usepackage{pgfplots}            % gráficos de datos dentro de LaTeX
\pgfplotsset{compat=1.18}
\usetikzlibrary{babel}            % evita conflictos de tikz con babel-spanish

% ---------- Citas y enlaces (hyperref siempre al final) ----------
\usepackage{cite}
\usepackage{url}
\usepackage{hyperref}

\title{\textbf{Diseño CAD/CAE de un sistema semiautomatizado para el recargue de soldadura en alojamientos de pasadores de maquinaria pesada}}
\author{Alba Rodríguez, Mauri Yeison \and Sandi Vargas, Julio Constantino}
\date{\today}

\begin{document}

\maketitle
\input{00-resumen/resumen.txt}

\newpage
\tableofcontents
\newpage

\input{01-introduccion/introduccion.txt}
\input{02-antecedentes/antecedentes.txt}
\input{03-problema/problema.txt}
\input{04-objetivos/objetivos.txt}
\input{05-necesidades-especificaciones/necesidades-especificaciones.txt}
\input{06-descomposicion-funcional/descomposicion-funcional.txt}
\input{07-busqueda-externa/busqueda-externa.txt}
\input{08-busqueda-interna/busqueda-interna.txt}
\input{09-exploracion-sistematica/exploracion-sistematica.txt}
\input{10-filtrado-conceptos/filtrado-conceptos.txt}
\input{11-evaluacion-seleccion/evaluacion-seleccion.txt}
\input{12-diseno-cad/diseno-cad.txt}
\input{13-analisis-fea/analisis-fea.txt}
\input{14-planos/planos.txt}
\input{15-conclusiones-recomendaciones/conclusiones-recomendaciones.txt}

\newpage
% Una sola lista de referencias en estilo IEEE, armada con los
% referencias.bib de todas las secciones.
\renewcommand{\refname}{Referencias}
\bibliographystyle{IEEEtran}
\bibliography{%
  00-resumen/referencias,%
  01-introduccion/referencias,%
  02-antecedentes/referencias,%
  03-problema/referencias,%
  04-objetivos/referencias,%
  05-necesidades-especificaciones/referencias,%
  06-descomposicion-funcional/referencias,%
  07-busqueda-externa/referencias,%
  08-busqueda-interna/referencias,%
  09-exploracion-sistematica/referencias,%
  10-filtrado-conceptos/referencias,%
  11-evaluacion-seleccion/referencias,%
  12-diseno-cad/referencias,%
  13-analisis-fea/referencias,%
  14-planos/referencias,%
  15-conclusiones-recomendaciones/referencias}

\end{document}
